% appendix_formula.tex: The formula with beliefs fixed, and its anchoring (2026-09-04).
\section{The Ratio Formula with Beliefs Fixed, and Its Anchoring}\label{appendix:formula}

This appendix derives the first-order estimates from the credit-bureau margins, the ratio formula with beliefs fixed, the anchoring of the term margin on the experiments, and the first-order run-off formula, reports what each returns on the data, and states why each is not the mapping the paper uses, a reason established in Appendix \ref{appendix:mcliquidity} or in Section \ref{sec:termmargin}. The estimates are reported so that the model-based estimates can be compared with them.

\subsection{The formula}
With an annuity contract, $m_s = m$ for $s \le T$, Proposition \ref{prop:sensT} writes each elasticity of the default probability as a weighted sum over the payment path. At decision month $s$ the borrower weighs the current payment, priced at the default threshold, against the remaining payments, priced at the repayers' marginal utility, $\lambda$ times the threshold value, discounted by $D(\cdot)$ and by survival in the loan $Q(\cdot)$:
\begin{equation}
W_s \;=\; 1 + \lambda \sum_{j=s+1}^{T} D(j-s)\,Q(j-s).
\label{eq:Ws}
\end{equation}
A penalty proportional to the remaining balance, $\kappa_B B_s$ with $B_s = m\,a(T-s; r)$ and $a(n; r) = (1-(1+r)^{-n})/r$ the annuity factor at the note rate $r$, adds a second channel to each elasticity. Writing $k = \kappa_B/\kappa$ for the ratio of the penalty scale to the payment scale, the ratio of the two elasticities at decision month $s$ is
\begin{equation}
R_s(\rho, k) \;=\; \frac{T\,\lambda\,D(T-s)\,Q(T-s) \;-\; k\,T\,(1+r)^{-(T-s)}}{W_s \;-\; k\,a(T-s; r)},
\label{eq:Rs}
\end{equation}
because the marginal payment added by a longer term falls at the end of the contract, and the remaining balance rises with the term at the rate $T(1+r)^{-(T-s)}/a(T-s;r)$; at $k = 0$ and $s = 1$, \eqref{eq:Rs} is the date-1 case of equation \eqref{eq:window} of the main text. The ratio for default within $K$ months aggregates decision months with the default-timing weights $\omega_s$, the hazard at $s$ times survival to $s$, normalized over $s \le K$:
\begin{equation}
R(K) \;\equiv\; \frac{\varepsilon_T(K)}{\varepsilon_m(K)} \;=\; \sum_{s \le K} \omega_s\, R_s(\rho, k).
\label{eq:RK}
\end{equation}
The penalty scale is reported as its share of the payment channel at the first decision month, $\phi = k\,a(T-1;r)/W_1 \in [0, 1)$, which makes the parameter space the same for every contract length and exit hazard. The inputs and their sources are: $\lambda = 0.2$ as a constant, the level the pair of experiments in Section \ref{sec:quasi} identifies, with a band from $0.1$ to $0.35$ (the anchored fit predates the profile of equation \eqref{eq:lambda_p} and is recorded as run); $Q(n) = e^{-h n/12}$ with $h$ the group's own exit hazard from payoff before maturity, corrected for exposure; $r = 0.08$ with a band from $0.05$ to $0.12$; $T$ the group's mean contract length; $\omega_s$ from the group's own default timing; and $D(n) = e^{-\rho n/12}$, the object of the fit. The asymmetry between the current payment and the later ones in \eqref{eq:Ws} is what the closed form $x/(e^x-1)$ omits; on the model's own simulations that omission overstates the ratio by a factor between 2.5 and 3 (Appendix \ref{appendix:mcliquidity}), and it is inside the formula here.

\subsection{What the profile identifies on its own}
Both channels in \eqref{eq:Rs} rise with the window: the preference channel because the marginal payment draws nearer at the rate $\rho + h$, the penalty channel because the balance amortizes in the denominator. Across the windows available, six to forty-eight months, a higher $\rho$ with a lower $\phi$ traces nearly the same curve as a lower $\rho$ with a higher $\phi$; the root mean squared distance between the best-matching curves at $\rho = 0$ and $\rho = 0.3$ is $0.04$, below the precision of any group. Fitted to the within-borrower profile, the bridged population profile, and every group profile, the formula reproduces each ($\chi^2$ below one on two degrees of freedom) while the confidence set for $\rho$ spans the whole grid. Contract-length bins do not resolve it: with an exit hazard near $0.22$ per year the weight on the marginal payment falls with the term through $Q$ rather than through $D$, so the gradient of $R$ in $T$ is nearly the same at every $\rho$. The profile therefore identifies the difference between the two channels: given $\rho$ it gives the penalty share, and given the penalty share it gives $\rho$.

\subsection{The anchoring}
Two measured objects supply what the profile lacks. The first moves the penalty channel and nothing else. Loans below a state's deficiency threshold cannot be pursued for the balance after repossession, so their default cost scales less with the remaining balance while their payment stream is unchanged; the difference in the ratio between these loans and loans of the same size in other states, $\Delta R(K) = R_{\text{no pursuit}}(K) - R_{\text{pursuit}}(K)$, is a function of $(\phi_{\text{pursuit}}, \phi_{\text{no pursuit}})$ at any $\rho$ (Section \ref{sec:termmargin}). The second anchors $\rho$: the rescheduling-against-write-down pair in Section \ref{sec:quasi}, read at a single decision date as the anchoring was run, fits best at $\rho^\ast = 0$ and accepts rates below $0.03$ at the five percent level and below $0.08$ at one percent over a four-year horizon for a distressed population; aggregated over the pair's own outcome window the set extends to $1.2$, which is why the paper no longer uses the pair as an anchor.

In the anchor group, the pooled sample or the bottom score tercile, the population closest to that experiment's, $\rho$ is set to $\rho^\ast$ and $(\phi_A, \phi_{\text{np},A})$ are fitted to the profile and the contrast. Their difference is the deficiency component of the balance channel, and $\psi = \phi_A/(\phi_A - \phi_{\text{np},A})$ is the ratio of the whole balance channel, whatever scales with the balance, to its deficiency part. In every other group $\psi$ is imposed, $\phi_{\text{np},g} = \phi_g(1 - 1/\psi)$, and $(\rho_g, \phi_g)$ are fitted to the group's own profile and contrast with the group's own exit hazard, timing weights and contract length. An alternative anchoring imposes the no-pursuit share itself, $\phi_{\text{np},g} = \phi_{\text{np},A}$; agreement between the two is reported as the check on the assumption. Groups too thin for a contrast inherit $\phi$ from their parent score tercile and fit $\rho$ alone. Uncertainty combines the profile confidence set given the anchoring with the band over $\rho^\ast$ from zero to the date-1 one percent bound of $0.08$ and over $\lambda \in \{0.1, 0.2, 0.35\}$.

The construction was to rest on three facts and one assumption. The discount rate of the anchor is identified by a rule-based experiment on a population the anchor group resembles. The payment margin, the object the experiments manipulate, is a stable parameter in the bureau data and is reproduced across the population by the specification of Appendix \ref{appendix:bridge}. The recourse contrast is the one bureau moment that moves a single channel. The assumption is that the balance channel scales with its deficiency component; it is imposed in two ways, and rankings of groups that survive both do not depend on it. What the construction does not provide is point identification of a discount rate for each group from the bureau alone: the population-scale statement is a discount rate across people and time given the population-level anchor.

\subsection{Recovery on simulated data, and the failure of the term-margin anchoring}
Two checks were run. On profiles generated from \eqref{eq:RK} with known parameters by group and noise at the precision of the sample, the anchored fit recovers the ordering of $\rho$ across groups and covers the true value in five of six groups. On the model of Appendix \ref{appendix:mcliquidity} itself, with the balance penalty at its full value for pursued loans and at forty percent of it for loans that cannot be pursued, the anchored fit was applied to seven simulated groups with known $(\rho, \sigma)$: it recovers the groups that share the anchor's penalty level ($0.30 \to 0.41$, $0.17 \to 0.19$, $0.03 \to 0.00$) and fails for the groups at twice the penalty ($0.08 \to 1.85$, $0.17 \to 1.60$, $0.30 \to 1.80$), with the anchor's own fit rejected ($\chi^2$ of 575 on twelve moments). The first-order formula is not a faithful reduced form of the model's window profile. On the covariate-matched extract, moreover, the recourse contrast has 303 no-pursuit defaults and no power by group, and the population profile exceeds one above the fifth score decile, which the preference channel cannot produce at any $\rho$. The term margin is therefore not used to identify a discount rate; its profile is reported as a finding in Section \ref{sec:termmargin}, and the designs applied to it are listed in Appendix \ref{appendix:strategies}.

\subsection{The payment margin over the life of the loan}\label{appendix:runoff}
Under the first-order formula $\rho$ would be identified from the payment margin itself, over the life of the loan. At loan age $s$ the incentive to default is $m W_s - \theta C_s$, with $W_s$ from \eqref{eq:Ws} evaluated over the $T - s$ remaining payments and $C_s$ the penalty base, measured by the charged-off amount. Under a logistic latent with a scale that does not change with loan age, the elasticity of the default hazard at age $s$ to $\log m$ is proportional to $W_s - \theta (C_s/m)\,\varepsilon_m$, where $\varepsilon_m = 1.04$ is the measured elasticity of the charged-off amount to the payment. $W_s$ runs off as the contract shortens, at a rate set by $\rho + h$ and scaled by $\lambda$; the penalty base runs off with the balance. Under that formula the shape of the run-off, not its level, would identify $\rho$: when the far payments have little weight, $W_s$ stays flat until the last two years and then drops, and when they have much, it declines from the start. Remark \ref{rem:option} and Appendix \ref{appendix:mcliquidity} establish that the stopping problem calibrated to the sample's default rate, with i.i.d.\ or with persistent income, produces the same shape at every rate, at a payment elasticity ten times the data's (Appendix \ref{appendix:persistent}).

The run-off is measured on the auto panel by a linear probability model of default within each loan-age bin (0--6, 6--12, 12--18, 18--24, 24--36, 36--48, 48--60 months), among loans that survived to the bin, on $\log m$ and $\log T$ with the fixed effects of Appendix \ref{appendix:bridge}, by group; the elasticity is the coefficient over the bin's hazard. The fit is
\begin{equation}
\varepsilon_m(s) \;=\; A\,\big[\,W_s(\rho, \lambda, h) - \theta\,(C_s/m)\,\varepsilon_m\,\big],
\label{eq:runoff}
\end{equation}
by weighted least squares on a $(\rho, \theta)$ grid with the scale $A$ in closed form, $\lambda = 0.2$, $h$ the group's own exit hazard, $T$ its mean contract length and $C_s/m$ from the charged-off share by loan age scaled to the group's level. The confidence set for $\rho$ is the $\chi^2(1)$ cut of the profile. Figure \ref{fig:runoff_profiles} in the main text shows the measured run-off for the pooled sample, the two score terciles and three income quartiles against the model's range: flat for three years and then falling, with the same shape in every group.



\subsection{Results}
The measured run-off and the model's reproduction of it are reported in Sections \ref{sec:mainresults} and \ref{sec:heterogeneity}. Fitted with the first-order formula \eqref{eq:runoff} at $\lambda = 0.2$ and the panel's exit hazard, the pooled profile on 61- to 72-month contracts returns an annual rate of $0.56$ with a set from $0.40$ to $0.84$, and the group profiles return rates from $0.1$ to $1.3$; these numbers are not rates, because Appendix \ref{appendix:mcliquidity} shows the formula returns rates near $3$ on the model's own simulations at true rates of $0.03$ to $0.10$. The paper matches the profile to the model's own run-off, and every group's set spans the grid.

Three caveats accompany the result. The proportionality of the elasticity to the incentive rests on a latent scale that does not change with loan age. First-payment defaults, which do not respond to the payment, lower the first bin and cost fit in the pooled sample ($\chi^2$ of 8.8 on four degrees of freedom). The exit hazard computed from monthly sightings is a lower bound, and the outcome is charge-off; the monthly one-percent panel with closure dates and first 90-day delinquency is the check on both.
